\subsection{7.1  Simulation Design}

The simulation advances in discrete ticks, each representing one operational day. Aria processes a fixed volume of synthetic customer interactions per tick, drawn from a weighted distribution across four task categories: order status (40\%), return requests (30\%), billing disputes (20\%), and product inquiries (10\%). Each interaction produces a PerformanceSnapshot recording task completion outcome, simulated CSAT score, policy citation accuracy, tool invocation correctness, and hallucination occurrence. Snapshots are evaluated against Aria's behavioral contract at the end of each tick; threshold violations trigger APIP instantiation. Token and latency features are extracted from simulated LLM call spans, consistent with the 16-feature trace-based failure taxonomy for agentic systems [31].

The DegradationEngine injects each of the four failure modes according to a parameterized schedule. Each mode degrades a distinct subset of metrics following a characteristic trajectory: Behavioral Drift ramps slowly then accelerates (simulating stale RAG); Input Brittleness drops performance sharply on a single task category while leaving overall metrics superficially intact; Alignment Creep produces a slow linear drift detectable only through the relative drift trigger; Tool Misuse spikes the tool error rate suddenly while leaving CSAT and task completion initially unaffected. Each scenario was run 30 times with randomized noise to produce distributions over outcomes.

Three operational cost metrics are tracked throughout the simulation in addition to the quality metrics in the behavioral contract. Token consumption measures the cumulative prompt and completion tokens consumed per tick, which increases under degradation as agents produce longer, more uncertain responses and require more retrieval context. Response latency (p95, in milliseconds) tracks the 95th percentile response time, which increases as hallucination and context retrieval load grows. Customer satisfaction loss (CSAT delta) measures the gap between the agent's current CSAT and its deployment baseline, providing a direct proxy for business impact. These three cost metrics ground the APIP's value proposition: earlier detection and correct-tier intervention reduces all three.

